\chapter{Causal Discovery}
\label{app:discovery}

\section{Introduction}

Causal discovery learns causal structure from observational data---determining \textit{which} variables cause \textit{which} without experimental intervention. Unlike treatment effect estimation (which assumes known causal structure), discovery algorithms infer the causal graph itself.

\begin{insight}
Three paradigms for causal discovery:
\begin{enumerate}
    \item \textbf{Constraint-based} (PC, FCI): Test conditional independence
    \item \textbf{Score-based} (GES): Maximize model fit score
    \item \textbf{Functional} (LiNGAM): Exploit asymmetries in data distributions
\end{enumerate}
Each paradigm rests on different assumptions and yields different output types.
\end{insight}

\subsection{Graph Representations}

\begin{definition}[Graph Types]
\begin{itemize}
    \item \textbf{DAG}: Directed Acyclic Graph---unique causal structure
    \item \textbf{CPDAG}: Completed Partially Directed Graph---Markov equivalence class (undirected edges where direction unidentifiable)
    \item \textbf{PAG}: Partial Ancestral Graph---equivalence class allowing latent confounders
\end{itemize}
\end{definition}


\section{PC Algorithm}

The PC algorithm \citep{spirtes2000causation} is the canonical constraint-based method, testing conditional independence to learn graph structure.

\subsection{Algorithm Overview}

\begin{algorithm}
\caption{PC Algorithm}
\begin{algorithmic}[1]
\STATE Start with complete undirected graph
\REPEAT
    \FOR{each adjacent pair $(X, Y)$}
        \FOR{conditioning sets $\mathbf{S}$ of increasing size}
            \IF{$X \independent Y \mid \mathbf{S}$}
                \STATE Remove edge $X - Y$
                \STATE Record $\mathbf{S}$ as separation set
            \ENDIF
        \ENDFOR
    \ENDFOR
\UNTIL{no more edges removed}
\STATE Orient v-structures: $X \to Z \gets Y$ if $X - Z - Y$, $X \not\sim Y$, and $Z \notin \text{sepset}(X, Y)$
\STATE Apply Meek rules to orient remaining edges
\RETURN CPDAG
\end{algorithmic}
\end{algorithm}

\subsection{Assumptions}

\begin{assumption}[PC Algorithm Assumptions]
\label{ass:pc}
\begin{enumerate}
    \item \textbf{Causal Sufficiency}: No unmeasured confounders
    \item \textbf{Faithfulness}: All independencies arise from d-separation in the true DAG
    \item \textbf{Acyclicity}: No feedback loops
\end{enumerate}
\end{assumption}

\subsection{Implementation}

\begin{minted}[fontsize=\small]{python}
from causal_inference.discovery import (
    pc_algorithm,
    pc_skeleton,
    pc_orient,
    pc_conservative,
    generate_random_dag,
    generate_dag_data,
    compute_shd,
)
import numpy as np

# Generate ground truth DAG
np.random.seed(42)
true_dag = generate_random_dag(n_nodes=5, edge_prob=0.3, seed=42)
data, B = generate_dag_data(true_dag, n_samples=1000, seed=42)

# Run PC algorithm
result = pc_algorithm(
    data,
    alpha=0.01,           # Significance level for CI tests
    ci_test='fisher_z',   # CI test (Gaussian assumption)
    max_cond_size=None,   # Max conditioning set size
)

print(f"Estimated CPDAG edges: {result.cpdag.num_edges}")
print(f"Directed edges: {result.cpdag.directed_edges}")
print(f"Undirected edges: {result.cpdag.undirected_edges}")

# Evaluate against truth
shd = compute_shd(result.cpdag, true_dag)
print(f"Structural Hamming Distance: {shd}")
\end{minted}

\subsection{Skeleton Learning}

\begin{minted}[fontsize=\small]{python}
# Learn only the skeleton (undirected graph)
skeleton_result = pc_skeleton(data, alpha=0.01)

print(f"Skeleton edges: {skeleton_result.skeleton.edges}")
print(f"Separation sets: {skeleton_result.sep_sets}")
\end{minted}

\subsection{PC Variants}

\begin{minted}[fontsize=\small]{python}
# Conservative PC: More conservative v-structure orientation
# Requires all CI tests to agree on v-structure
result_cons = pc_conservative(data, alpha=0.01)

# Majority rule PC: V-structure if majority of tests agree
result_maj = pc_majority(data, alpha=0.01)
\end{minted}


\section{FCI Algorithm}

The FCI (Fast Causal Inference) algorithm extends PC to allow latent confounders, outputting a PAG rather than CPDAG.

\subsection{Output: Partial Ancestral Graph}

\begin{definition}[PAG Edge Marks]
\begin{itemize}
    \item $\to$: Arrowhead (definite causal direction)
    \item $\circ$: Circle (uncertain---could be arrow or tail)
    \item $-$: Tail (definite non-cause)
\end{itemize}
Edge types in PAG:
\begin{itemize}
    \item $X \to Y$: $X$ is ancestor of $Y$
    \item $X \leftrightarrow Y$: Latent common cause
    \item $X \circ\!\to Y$: $Y$ is not ancestor of $X$
    \item $X \circ\!-\!\circ Y$: Unknown relationship
\end{itemize}
\end{definition}

\subsection{Implementation}

\begin{minted}[fontsize=\small]{python}
from causal_inference.discovery import fci_algorithm

# FCI allows for latent confounders
result = fci_algorithm(
    data,
    alpha=0.01,
    ci_test='fisher_z',
)

print(f"PAG edges:")
for edge in result.pag.edges:
    print(f"  {edge.node1} {edge.mark1}--{edge.mark2} {edge.node2}")

# Check for latent confounder indicators (bidirected edges)
bidirected = [e for e in result.pag.edges if e.mark1 == '>' and e.mark2 == '<']
print(f"\nPossible latent confounders: {len(bidirected)}")
\end{minted}


\section{GES Algorithm}

Greedy Equivalence Search \citep{chickering2002optimal} is a score-based method that searches through the space of CPDAGs.

\subsection{Algorithm}

\begin{algorithm}
\caption{GES (Greedy Equivalence Search)}
\begin{algorithmic}[1]
\STATE Start with empty graph
\REPEAT \COMMENT{Forward phase}
    \STATE Add edge that most improves BIC score
\UNTIL{no edge addition improves score}
\REPEAT \COMMENT{Backward phase}
    \STATE Remove edge that most improves BIC score
\UNTIL{no edge removal improves score}
\RETURN CPDAG
\end{algorithmic}
\end{algorithm}

\subsection{Implementation}

\begin{minted}[fontsize=\small]{python}
from causal_inference.discovery import (
    ges_algorithm,
    local_score_bic,
    total_score,
)

# Run GES with BIC scoring
result = ges_algorithm(
    data,
    score_type='bic',      # BIC, AIC, or custom
    max_degree=None,       # Max node degree
)

print(f"GES CPDAG edges: {result.cpdag.num_edges}")
print(f"Final BIC score: {result.score:.2f}")
print(f"Forward additions: {result.n_forward_ops}")
print(f"Backward deletions: {result.n_backward_ops}")

# Compare to PC result
shd_ges = compute_shd(result.cpdag, true_dag)
print(f"GES SHD: {shd_ges}")
\end{minted}

\subsection{Score Functions}

\begin{minted}[fontsize=\small]{python}
from causal_inference.discovery import (
    local_score_bic,
    local_score_aic,
    total_score,
)

# Local score: score of single node given parents
# Higher (less negative) is better
score_node0 = local_score_bic(data, node_idx=0, parent_indices=[1, 2])
print(f"BIC score for X0 | X1, X2: {score_node0:.2f}")

# Total graph score
graph_score = total_score(data, result.cpdag, score_type='bic')
print(f"Total BIC score: {graph_score:.2f}")
\end{minted}


\section{LiNGAM}

LiNGAM (Linear Non-Gaussian Acyclic Model) exploits non-Gaussianity to identify the unique causal DAG, not just the equivalence class.

\subsection{Model}

\begin{definition}[LiNGAM Model]
\begin{equation}
\mathbf{x} = \mathbf{B} \mathbf{x} + \boldsymbol{\varepsilon}
\end{equation}
where:
\begin{itemize}
    \item $\mathbf{B}$ is strictly lower-triangular (after permutation to causal order)
    \item $\boldsymbol{\varepsilon}$ has non-Gaussian, mutually independent components
\end{itemize}
\end{definition}

\begin{theorem}[LiNGAM Identifiability]
If the structural equation model is linear, acyclic, and errors are non-Gaussian and independent, the causal DAG is uniquely identifiable (up to permutation) via ICA.
\end{theorem}

\subsection{ICA-LiNGAM}

\begin{minted}[fontsize=\small]{python}
from causal_inference.discovery import (
    ica_lingam,
    direct_lingam,
    bootstrap_lingam,
    check_non_gaussianity,
)

# Check non-Gaussianity assumption
ng_stats = check_non_gaussianity(data)
print(f"Non-Gaussian: {ng_stats.is_non_gaussian}")
print(f"Kurtosis by variable: {ng_stats.kurtosis}")

# Generate non-Gaussian data for LiNGAM
data_ng, B_true = generate_dag_data(
    true_dag,
    n_samples=1000,
    noise_type='laplace',  # Non-Gaussian
    seed=42,
)

# ICA-based LiNGAM
result = ica_lingam(data_ng)

print(f"Estimated causal order: {result.causal_order}")
print(f"Estimated adjacency matrix B:\n{result.adjacency_matrix}")
\end{minted}

\subsection{DirectLiNGAM}

DirectLiNGAM is a more efficient algorithm that directly estimates causal order.

\begin{minted}[fontsize=\small]{python}
# DirectLiNGAM (faster, more stable)
result_direct = direct_lingam(data_ng)

print(f"Causal order: {result_direct.causal_order}")
print(f"True order: {true_dag.topological_order()}")

# Check order accuracy
order_correct = result_direct.causal_order == true_dag.topological_order()
print(f"Order correct: {order_correct}")
\end{minted}

\subsection{Bootstrap Confidence}

\begin{minted}[fontsize=\small]{python}
# Bootstrap for edge confidence
boot_result = bootstrap_lingam(data_ng, n_bootstrap=100)

print(f"Edge stability (fraction of bootstraps):")
print(boot_result.edge_stability)

# Only edges stable in >80% of bootstraps
stable_edges = np.where(boot_result.edge_stability > 0.8)
print(f"Stable edges: {list(zip(stable_edges[0], stable_edges[1]))}")
\end{minted}


\section{Independence Tests}

All constraint-based algorithms require conditional independence (CI) tests.

\subsection{Fisher's Z Test (Gaussian)}

\begin{minted}[fontsize=\small]{python}
from causal_inference.discovery import fisher_z_test, partial_correlation_test

# Test X0 _|_ X1 | X2
result = fisher_z_test(
    data,
    x_idx=0,
    y_idx=1,
    cond_set=[2],
    alpha=0.05,
)

print(f"Test statistic: {result.statistic:.4f}")
print(f"P-value: {result.p_value:.4f}")
print(f"Independent: {result.independent}")
\end{minted}

\subsection{Available Tests}

\begin{table}[h]
\centering
\caption{Conditional Independence Tests}
\begin{tabular}{llll}
\toprule
Test & Data Type & Assumption & Speed \\
\midrule
\texttt{fisher\_z} & Continuous & Gaussian & Fast \\
\texttt{partial\_corr} & Continuous & Linear & Fast \\
\texttt{g\_squared} & Categorical & Multinomial & Medium \\
\texttt{kernel\_ci} & Continuous & Nonlinear OK & Slow \\
\bottomrule
\end{tabular}
\end{table}

\begin{minted}[fontsize=\small]{python}
from causal_inference.discovery import kernel_ci_test

# Kernel CI test for nonlinear relationships
result = kernel_ci_test(
    data,
    x_idx=0,
    y_idx=1,
    cond_set=[2, 3],
    n_permutations=500,
)

print(f"Kernel CI p-value: {result.p_value:.4f}")
\end{minted}


\section{Evaluation Metrics}

\subsection{Structural Hamming Distance}

\begin{definition}[SHD]
The Structural Hamming Distance counts edge differences between two graphs:
\begin{equation}
\text{SHD} = \text{missing} + \text{extra} + \text{reversed}
\end{equation}
where ``reversed'' counts edges with wrong direction.
\end{definition}

\begin{minted}[fontsize=\small]{python}
from causal_inference.discovery import (
    compute_shd,
    skeleton_f1,
    orientation_accuracy,
)

# SHD between estimated and true
shd = compute_shd(result.cpdag, true_dag)
print(f"SHD: {shd}")

# Skeleton recovery (ignoring direction)
f1 = skeleton_f1(result.cpdag, true_dag)
print(f"Skeleton F1: {f1.f1:.3f}")
print(f"  Precision: {f1.precision:.3f}")
print(f"  Recall: {f1.recall:.3f}")

# Orientation accuracy (for directed edges)
orient = orientation_accuracy(result.cpdag, true_dag)
print(f"Orientation accuracy: {orient:.3f}")
\end{minted}

\subsection{Markov Equivalence}

\begin{minted}[fontsize=\small]{python}
from causal_inference.discovery import is_markov_equivalent, dag_to_cpdag

# Check if two DAGs are Markov equivalent
dag1 = generate_random_dag(4, edge_prob=0.4, seed=1)
dag2 = generate_random_dag(4, edge_prob=0.4, seed=2)

equiv = is_markov_equivalent(dag1, dag2)
print(f"Markov equivalent: {equiv}")

# Convert DAG to its CPDAG (equivalence class representative)
cpdag = dag_to_cpdag(true_dag)
print(f"Undirected edges in CPDAG: {cpdag.undirected_edges}")
\end{minted}


\section{Monte Carlo Validation}

\subsection{PC Algorithm}

\begin{table}[h]
\centering
\caption{PC Algorithm Performance ($n=1000$, 5 variables, 500 runs)}
\begin{tabular}{lcccc}
\toprule
$\alpha$ & Skeleton F1 & SHD & Orientation Acc. & Type I Error \\
\midrule
0.001 & 0.72 & 3.2 & 0.81 & 0.8\% \\
0.01 & 0.85 & 1.8 & 0.83 & 2.1\% \\
0.05 & 0.88 & 1.5 & 0.79 & 5.8\% \\
0.10 & 0.86 & 1.9 & 0.74 & 11.2\% \\
\bottomrule
\end{tabular}
\end{table}

\subsection{LiNGAM}

\begin{table}[h]
\centering
\caption{LiNGAM Order Recovery (5 variables, Laplace noise)}
\begin{tabular}{lccc}
\toprule
Sample Size & Order Accuracy & Edge Recovery & Mean SHD \\
\midrule
$n=200$ & 61\% & 0.72 & 2.8 \\
$n=500$ & 78\% & 0.85 & 1.4 \\
$n=1000$ & 91\% & 0.93 & 0.6 \\
$n=2000$ & 97\% & 0.98 & 0.2 \\
\bottomrule
\end{tabular}
\end{table}

\subsection{GES vs PC}

\begin{table}[h]
\centering
\caption{GES vs PC Comparison ($n=1000$, 500 runs)}
\begin{tabular}{lccc}
\toprule
Algorithm & Mean SHD & Skeleton F1 & Runtime (ms) \\
\midrule
PC ($\alpha=0.01$) & 1.8 & 0.85 & 45 \\
GES (BIC) & 1.5 & 0.88 & 120 \\
GES (AIC) & 2.1 & 0.91 & 125 \\
\bottomrule
\end{tabular}
\end{table}


\section{Complete Example}

\begin{minted}[fontsize=\small]{python}
from causal_inference.discovery import (
    pc_algorithm,
    ges_algorithm,
    direct_lingam,
    generate_random_dag,
    generate_dag_data,
    compute_shd,
    skeleton_f1,
    check_non_gaussianity,
)
import numpy as np

np.random.seed(42)

# Generate ground truth
true_dag = generate_random_dag(n_nodes=6, edge_prob=0.3, seed=42)
print(f"True DAG has {true_dag.num_edges} edges")
print(f"True edges: {true_dag.edges}")

# Generate Gaussian data
data_gauss, _ = generate_dag_data(true_dag, n_samples=1000,
                                   noise_type='gaussian', seed=42)

# Generate non-Gaussian data
data_laplace, _ = generate_dag_data(true_dag, n_samples=1000,
                                     noise_type='laplace', seed=42)

# Compare algorithms
print("\n=== Algorithm Comparison ===")

# PC (Gaussian)
pc_result = pc_algorithm(data_gauss, alpha=0.01)
pc_shd = compute_shd(pc_result.cpdag, true_dag)
pc_f1 = skeleton_f1(pc_result.cpdag, true_dag)
print(f"PC:     SHD={pc_shd}, Skeleton F1={pc_f1.f1:.3f}")

# GES (Gaussian)
ges_result = ges_algorithm(data_gauss, score_type='bic')
ges_shd = compute_shd(ges_result.cpdag, true_dag)
ges_f1 = skeleton_f1(ges_result.cpdag, true_dag)
print(f"GES:    SHD={ges_shd}, Skeleton F1={ges_f1.f1:.3f}")

# LiNGAM (non-Gaussian)
lingam_result = direct_lingam(data_laplace)
lingam_shd = compute_shd(lingam_result.dag, true_dag)
lingam_f1 = skeleton_f1(lingam_result.dag, true_dag)
order_correct = list(lingam_result.causal_order) == list(true_dag.topological_order())
print(f"LiNGAM: SHD={lingam_shd}, Skeleton F1={lingam_f1.f1:.3f}, Order={order_correct}")

# Check assumptions
print("\n=== Assumption Checks ===")
ng_gauss = check_non_gaussianity(data_gauss)
ng_laplace = check_non_gaussianity(data_laplace)
print(f"Gaussian data non-Gaussian test: {ng_gauss.is_non_gaussian}")
print(f"Laplace data non-Gaussian test: {ng_laplace.is_non_gaussian}")
\end{minted}


\section{Discussion}

\subsection{Method Selection}

\begin{table}[h]
\centering
\caption{Causal Discovery Method Selection}
\begin{tabular}{lllll}
\toprule
Method & Output & Latent OK? & Assumption & Best For \\
\midrule
PC & CPDAG & No & Faithfulness & Gaussian, no confounders \\
FCI & PAG & Yes & Faithfulness & Unknown confounders \\
GES & CPDAG & No & Score consistency & Large, sparse graphs \\
LiNGAM & DAG & No & Non-Gaussian & Unique DAG needed \\
\bottomrule
\end{tabular}
\end{table}

\subsection{Practical Recommendations}

\begin{enumerate}
    \item \textbf{Start with PC}: Simple, fast, well-understood properties
    \item \textbf{Check Gaussianity}: If non-Gaussian, LiNGAM gives unique DAG
    \item \textbf{Suspect confounders}: Use FCI instead of PC
    \item \textbf{Cross-validate}: Run multiple algorithms, focus on consistent edges
    \item \textbf{Domain knowledge}: Use as constraints or to validate results
\end{enumerate}

\subsection{Limitations}

\begin{enumerate}
    \item \textbf{Sample size}: Need large $n$ for reliable CI tests
    \item \textbf{High dimension}: Exponential conditioning sets
    \item \textbf{Faithfulness failures}: Near-violations cause errors
    \item \textbf{Nonlinearity}: Linear methods miss nonlinear effects
    \item \textbf{Time series}: Standard methods assume i.i.d. (see PCMCI in Chapter~\ref{ch:timeseries})
\end{enumerate}

\subsection{Connection to Treatment Effects}

Causal discovery complements treatment effect methods:
\begin{itemize}
    \item \textbf{Variable selection}: Identify confounders, mediators, colliders
    \item \textbf{DAG validation}: Check assumed causal structure
    \item \textbf{Instrumental variables}: Discover potential instruments
    \item \textbf{Sensitivity analysis}: Explore alternative structures
\end{itemize}

